% Insert after Theorem SA21-A14 and its proof.
% AI-produced by GPT-6 Astra (OpenAI), ChatGPT Pro mode, 19 November 2026.
\begin{corollary}[Proved: the exact atom class for interval magasins]
\label{SA21-A14a}
For every poset $L$, the lieux of $M(L)$ are exactly the point symbols
$P_x$. All these lieux have atomic supports if and only if $L$ is an
antichain or $L$ has no isolated points. Thus every lieu has an atomic
support for every nonempty connected poset, including the singleton.
For $L=\varnothing$ the assertion is vacuous.
\end{corollary}
\begin{proof}
We use the manuscript's full-refinement definition of a lieu:
minimality for $\ll$, not merely for interior refinement $\rho$
(156-8, scan~15, \path{J156-8-02_transcription.tex:277}).
In the primed formulation it is minimality for both incidence and
$\rho$ (\path{J156-8-04_transcription.tex:282,296}).
Only $P_x$ refines $P_x$, since a factorization
$U\rho Z\le P_x$ forces $Z=P_x$ and then $U=P_x$.
Every interval $I_{ab}$ has the distinct incidence predecessor
$P_a$, which also refines it. Thus the lieux are exactly the points,
even when $a\lessdot b$ makes $I_{ab}$ minimal for $\rho$ alone.
In particular we make no unsupported divisibility claim for intervals
having no point strictly between their endpoints.

Theorem~\ref{SA21-A06} proves that every nonisolated $x$ has
$\operatorname{cl}\{P_x\}=\{P_x\}$, a nonzero atom.
For an isolated $x$, the symbol $P_x$ has no neighbour, so
$\operatorname{cl}\{P_x\}=M(L)$. If $L$ has a comparison $a<b$,
the singleton support of $P_a$ is a proper nonzero closed set,
and this top support is not an atom. If $L$ is a nonempty antichain,
there are no interval symbols and no disjunction edges; its support
lattice has just zero and one, so every lieu has that single nonzero
atom as support. These alternatives prove both directions of the
criterion. A connected comparability graph with more than one vertex
has no isolated point, and the singleton case is an antichain.
The empty case follows from the universal quantifier over no lieux.
\end{proof}
